\apendice{Documentación de usuario}

\section{Introducción}
La presente sección constituye el manual de usuario de la aplicación
\nombreProyecto, desarrollada como parte del Trabajo de Fin de Grado del grado
de Ingeniería Informática de la Universidad de Burgos. El sistema consiste en
una aplicación web diseñada para explorar y comparar los resultados obtenidos
del análisis temático sobre memorias y repositorios de Trabajos Fin de Grado
disponibles públicamente en GitHub.

El proyecto se estructura en tres módulos principales de forma transparente
para el usuario final:
\begin{itemize}
    \item \textbf{Ingesta (\textit{backend-ingestion}):} Módulo encargado de la extracción de datos desde la \acrshort{api} de GitHub y su almacenamiento estructurado en ficheros formato Parquet.
    \item \textbf{Procesamiento (\textit{backend-processing}):} Módulo analítico que procesa los textos extraídos, entrena los modelos de modelado de temas (\acrshort{lda}, BERTopic, Top2Vec, FASTopic) y evalúa los hiperparámetros óptimos.
    \item \textbf{Interfaz Web (\textit{frontend-web}):} Aplicación interactiva desarrollada en Streamlit~\cite{StreamlitDocsa} que permite la visualización y exploración de los resultados almacenados.
\end{itemize}

\section{Requisitos del sistema}
Para asegurar el correcto funcionamiento de la plataforma, el equipo del
usuario debe cumplir con una serie de requisitos dependiendo del modo de uso
elegido:

\begin{itemize}
    \item \textbf{Modo Visualización:} Únicamente se requiere disponer de un navegador web moderno (Google Chrome, Mozilla Firefox, Edge) para acceder a una instancia de la aplicación ya desplegada.
    \item \textbf{Modo Recomendado (Docker):} Es necesario tener instalado \textit{Docker Engine} y \textit{Docker Compose} (o \textit{Docker Desktop} en entornos Windows/macOS). Esta es la forma más sencilla de ejecutar el proyecto, ya que utiliza imágenes precompiladas y no requiere configurar el entorno de programación de Python localmente.
    \item \textbf{Modo Desarrollo Local:} Requiere tener instalado Python 3.10, gestor de dependencias (\textit{pip} o \textit{Conda}), la herramienta \textit{make} y el sistema de control de versiones \textit{Git}.
\end{itemize}

\section{Instalación y despliegue}
El sistema se ha empaquetado utilizando contenedores
Docker~\cite{DockerManuals} alojados en \ACRfull{ghcr}, lo que permite un
despliegue rápido sin necesidad de compilar el código fuente.

Para ejecutar la aplicación completa en un equipo local mediante Docker, se
deben seguir los siguientes pasos:
\begin{enumerate}
    \item \textbf{Clonar el repositorio:} Obtener el código del proyecto desde el repositorio oficial mediante la línea de comandos:
          \begin{verbatim}
git clone https://github.com/zcc1001/tfg-topics
cd tfg-topics
    \end{verbatim}
    \item \textbf{Actualizar las imágenes:} Descargar las últimas versiones de los contenedores utilizando la herramienta \textit{Make}:
          \begin{verbatim}
make pull
    \end{verbatim}
    \item \textbf{Ejecutar el flujo de trabajo:} El repositorio incluye comandos automatizados. Para ejecutar la extracción de datos, el procesamiento de los modelos y lanzar la interfaz web de manera secuencial y desatendida, se debe ejecutar:
          \begin{verbatim}
make pipeline
    \end{verbatim}
    \item \textbf{Arrancar la interfaz aislada:} Si los datos ya han sido procesados previamente y se desea arrancar únicamente la interfaz visual, se utiliza el comando:
          \begin{verbatim}
make start
    \end{verbatim}
\end{enumerate}

Una vez finalizado el proceso de despliegue, la interfaz web estará disponible
accediendo a la \acrshort{url} \texttt{http://localhost:8501} desde cualquier
navegador web en la máquina anfitriona.

\section{Manual del usuario de la interfaz}
La aplicación web ofrece una interfaz gráfica intuitiva con un menú de
navegación situado en el panel lateral izquierdo. Este menú permite alternar
entre dos enfoques principales de análisis: el \textbf{Modo Académico} y el
\textbf{Modo Técnico} como se muestra en la
Figura~\ref{fig:manual_usuario/menu_lateral.png}.

\imagen{manual_usuario/menu_lateral.png}{Menu lateral de navegación}

\subsection{Modo Académico}
Este modo está orientado al análisis cualitativo y exploratorio de los
proyectos, y está pensado para docentes, miembros de tribunal o investigadores
que deseen comprender el contenido de los trabajos sin entrar en detalles
algorítmicos complejos.

Como se aprecia en la Figura~\ref{fig:manual_usuario/modo_academico.png}, la
interfaz proporciona:
\begin{itemize}
    \item \textbf{Filtros interactivos:} Permite afinar la búsqueda de proyectos utilizando controles interactivos. El usuario puede buscar por nombre de tutor (campo de texto opcional), seleccionar un rango de años de presentación y acotar por nota mediante un control deslizante de rango.
    \item \textbf{Exploración de proyectos:} Una vez aplicados los filtros, se muestra el listado de resultados. Cada trabajo se presenta en una tarjeta detallada que incluye el título, un resumen descriptivo y una ficha con la \textit{Información Académica} (Autor, Tutor, Centro, Curso, etc.).
    \item \textbf{Distribución de temáticas:} Muestra gráficos visuales y resúmenes sobre cuáles son las tecnologías, metodologías y temáticas más frecuentemente tratadas por el alumnado en los diferentes cursos académicos (ver Figura \ref{fig:manual_usuario/distribucion_academica.png}).
\end{itemize}

\begin{figure}[H]
    \centering
    \includegraphics[height=0.8\textheight]{manual_usuario/modo_academico.png}
    \caption{Vista Análisis Académico}
    \label{fig:manual_usuario/modo_academico.png}
\end{figure}
\vspace{-8pt}
\imagen{manual_usuario/distribucion_academica.png}{Vista Distribución Académica}

\subsection{Modo Técnico}
Orientado al análisis de los resultados y métricas obtenidas por los distintos
algoritmos de Procesamiento de Lenguaje Natural (\acrshort{nlp}). Al margen de
los controles globales situados en la parte inferior del menú lateral (como el
selector de \textbf{Idioma / Language} y el selector de \textbf{Vista}), una
vez activado el \textit{Análisis Técnico}, el menú se despliega ofreciendo una
estructura jerárquica dividida en tres bloques principales:

\imagen{manual_usuario/modo_tecnico.png}{Vista Modo Técnico}

\subsubsection{Análisis por modelo}
En esta sección el usuario puede explorar los resultados detallados de un
algoritmo y conjunto de datos específicos. Incluye las siguientes opciones de
navegación:
\begin{itemize}
    \item \textbf{Resumen ejecutivo:} (Figura~\ref{fig:manual_usuario/resumen_ejecutivo.png}) Encabeza este apartado ofreciendo métricas de alto nivel, tales como el número de temas detectados, la cantidad de documentos analizados y el valor de coherencia global del modelo. Además, destaca visualmente el tema con mayor peso o distribución.
    \item \textbf{Exploración de temas:} Proporciona una tabla interactiva (\textit{Tabla de temas}) que relaciona el identificador de cada tema con sus palabras clave más relevantes, y ofrece vistas alternativas como nubes de palabras (\textit{Wordcloud}, Figura \ref{fig:manual_usuario/wordcloud_topicos.png}).
    \item \textbf{Mapa intertópico:} Representación bidimensional para analizar la distancia semántica y relación entre los distintos temas descubiertos (Figura \ref{fig:manual_usuario/mapa_intertopico.png}).
\end{itemize}

\imagen{manual_usuario/resumen_ejecutivo.png}{Vista Resumen Ejecutivo - Análisis por modelo}

\imagen{manual_usuario/wordcloud_topicos.png}{Vista Nube de Palabras (Wordcloud) - Exploración de temas}

\imagen{manual_usuario/mapa_intertopico.png}{Vista Mapa Intertópico}

\subsubsection{Hiperparametrización}
Permite observar la evolución y los resultados del proceso de ajuste automático
de parámetros de los modelos, a través de las siguientes subsecciones:
\begin{itemize}
    \item \textbf{Evolución del score:} Muestra de forma general cómo ha mejorado la métrica objetivo a lo largo del proceso de optimización.
    \item \textbf{Parámetros vs rendimiento:} (Figura~\ref{fig:manual_usuario/evolucion_score.png}) Presenta una gráfica interactiva de dispersión (\textit{scatter plot}) que muestra el valor de coherencia obtenido en cada intento (\textit{trial}). Facilita la identificación del mejor modelo, marcado de forma destacada (por ejemplo, con una estrella roja indicando el \textit{Best trial}).
    \item \textbf{Mejores hiperparámetros:} (Figura~\ref{fig:manual_usuario/mejores_params.png}) Muestra la configuración exacta que maximizó el rendimiento del modelo durante el entrenamiento.
\end{itemize}

\imagen{manual_usuario/evolucion_score.png}{Vista Evolución del score - Hiperparametrización}
\imagen{manual_usuario/mejores_params.png}{Vista Mejores hiperparámetros - Hiperparametrización}
\subsubsection{Comparativa de modelos}
Aglutina las visiones que enfrentan los resultados de los cuatro modelos
empleados (\acrshort{lda}~\cite{blei2003lda},
BERTopic~\cite{QuickStartBERTopic},
Top2Vec~\cite{angelovTop2VecDistributedRepresentations2020} y
FASTopic~\cite{FastopicFASTopic}) bajo diferentes métricas cuantitativas.
Comprende las vistas de:
\begin{itemize}
    \item \textbf{Resumen comparativo} y \textbf{Ranking global:} Evalúan el nivel de diversidad temática y coherencia de manera conjunta para establecer qué algoritmo se comporta mejor (Figura \ref{fig:manual_usuario/ranking_modelos.png}).
    \item \textbf{Comparación de métricas:} Gráficos cruzados para contrastar los distintos aspectos de evaluación entre los algoritmos entrenados (Figura \ref{fig:manual_usuario/comparativa_modelos.png}).
\end{itemize}

\imagen{manual_usuario/ranking_modelos.png}{Vista Comparativa Modelo - Ranking modelos}
\imagen{manual_usuario/comparativa_modelos.png}{Vista Comparativa Modelos - Resumen comparativo}